\section{Results}
\label{sec:results}

\begin{table*}[t]
  \centering
  \caption{Main results, PaDiM vs PatchCore (WRN50, 224\,px, seed 0). Image F1 uses the threshold set on the normal validation split. \oracle{} selects its threshold with test labels and is shown for reference only. The better value of each pair is in bold.}
  \label{tab:main}
  \footnotesize
  \begin{tabular}{@{}l@{\hspace{4pt}}r@{\hspace{4pt}}r@{\hspace{4pt}}r@{\hspace{4pt}}r@{\hspace{4pt}}r@{\hspace{4pt}}r@{\hspace{4pt}}r@{\hspace{4pt}}r@{\hspace{4pt}}r@{\hspace{4pt}}r@{}}
\toprule
& \multicolumn{2}{c}{Image AUROC} & \multicolumn{2}{c}{Image F1 (val thr)} & \multicolumn{2}{c}{\oracle} & \multicolumn{2}{c}{Pixel AUROC} & \multicolumn{2}{c}{AU-PRO} \\
\cmidrule(lr){2-3}\cmidrule(lr){4-5}\cmidrule(lr){6-7}\cmidrule(lr){8-9}\cmidrule(lr){10-11}
Category & PaDiM & PatchCore & PaDiM & PatchCore & PaDiM & PatchCore & PaDiM & PatchCore & PaDiM & PatchCore \\
\midrule
bottle & 0.999 & \textbf{1.000} & 0.977 & \textbf{1.000} & 0.992 & \textbf{1.000} & 0.982 & \textbf{0.987} & 0.949 & \textbf{0.949} \\
capsule & 0.902 & \textbf{0.978} & 0.787 & \textbf{0.836} & 0.969 & \textbf{0.977} & 0.985 & \textbf{0.989} & 0.914 & \textbf{0.931} \\
screw & 0.834 & \textbf{0.970} & 0.222 & \textbf{0.884} & 0.905 & \textbf{0.963} & 0.984 & \textbf{0.993} & 0.935 & \textbf{0.963} \\
carpet & \textbf{0.999} & 0.987 & \textbf{0.983} & 0.935 & \textbf{0.994} & 0.977 & 0.989 & \textbf{0.990} & \textbf{0.960} & 0.950 \\
tile & 0.981 & \textbf{0.989} & 0.286 & \textbf{0.895} & 0.963 & \textbf{0.988} & 0.939 & \textbf{0.960} & 0.799 & \textbf{0.819} \\
leather & \textbf{1.000} & \textbf{1.000} & \textbf{0.995} & 0.963 & \textbf{1.000} & \textbf{1.000} & 0.989 & \textbf{0.993} & 0.971 & \textbf{0.972} \\
\midrule
Mean (obj.) & 0.912 & \textbf{0.983} & 0.662 & \textbf{0.907} & 0.955 & \textbf{0.980} & 0.983 & \textbf{0.990} & 0.933 & \textbf{0.948} \\
Mean (tex.) & \textbf{0.993} & 0.992 & 0.754 & \textbf{0.931} & 0.986 & \textbf{0.988} & 0.972 & \textbf{0.981} & 0.910 & \textbf{0.914} \\
Mean (all) & 0.953 & \textbf{0.987} & 0.708 & \textbf{0.919} & 0.971 & \textbf{0.984} & 0.978 & \textbf{0.985} & 0.921 & \textbf{0.931} \\
\bottomrule
\end{tabular}

\end{table*}

\begin{figure*}[t]
  \centering
  \includegraphics[width=\textwidth]{fig_main_dumbbell.pdf}
  \caption{Per-category comparison of PaDiM (circles) and PatchCore (squares): (a)~image AUROC, (b)~image F1 at the normal-validation threshold, (c)~AU-PRO. A segment joins the two methods for each category. Panels (a) and (c) share the range $[0.75,1]$ and panel (b) spans $[0,1]$. Values are single-seed estimates, so no uncertainty is shown. PaDiM falls behind on screw and capsule, the two methods tie on the textures, and PaDiM's F1 collapses on screw and tile despite usable AUROC.}
  \label{fig:dumbbell}
\end{figure*}

\subsection{Overall comparison}
Table~\ref{tab:main} and Fig.~\ref{fig:dumbbell} summarise the results.
PatchCore is better on every six-category mean: image AUROC \pcImgAurocAll{} vs \pdImgAurocAll{}, image F1 \pcImgFoneAll{} vs \pdImgFoneAll{}, pixel AUROC \pcPixAurocAll{} vs \pdPixAurocAll{} and AU-PRO \pcAuproAll{} vs \pdAuproAll{}.
It is also cheaper (Table~\ref{tab:cost}).
PatchCore fits about $1.7\times$ faster, runs at \pcInferAll{} vs \pdInferAll{}\,ms per image, and uses less peak GPU memory.
PaDiM's peak GPU figure understates its real footprint, because its 3.8\,GB of inverse covariances per category sit in host memory.

\begin{table}[t]
  \centering
  \caption{Cost, six-category means on an RTX~5080. Infer excludes and e2e includes data loading.}
  \label{tab:cost}
  \footnotesize
  \begin{tabular}{@{}lrrrr@{}}
\toprule
Method & \shortstack{Fit\\(s)} & \shortstack{Infer\\(ms/img)} & \shortstack{E2E\\(ms/img)} & \shortstack{GPU\\(GiB)} \\
\midrule
PaDiM & 6.15 & 3.58 & 10.10 & 1.72 \\
PatchCore & \textbf{3.63} & \textbf{2.78} & \textbf{9.31} & \textbf{1.50} \\
\bottomrule
\end{tabular}

\end{table}

\subsection{Objects vs textures}
\textbf{Textures are a tie.}
Mean texture image AUROC is \pdImgAurocTex{} for PaDiM and \pcImgAurocTex{} for PatchCore, and PaDiM is even slightly better on carpet (\pdImgAurocCarpet{} vs \pcImgAurocCarpet{}).
In a stationary texture every position sees the same statistics, so position-specific Gaussians lose nothing.

\textbf{The gap is in the objects.}
Mean object image AUROC is \pdImgAurocObj{} vs \pcImgAurocObj{}, driven by screw (\pdImgAurocScrew{} vs \pcImgAurocScrew{}) and capsule (\pdImgAurocCapsule{} vs \pcImgAurocCapsule{}).
This confirms the prediction of Section~\ref{sec:related}, although the two categories fail for different reasons.

\emph{Screw: pose.}
Screws are poorly aligned (align 0.579).
A fixed grid position therefore sees thread in one image, shank in another and background in a third, and its single Gaussian must stretch to cover all of them.
As a result, a damaged thread falls ``within normal variation'' at that position.
PatchCore asks only whether the damaged patch resembles \emph{any} normal patch, so pose does not matter.
PaDiM's image recall on screw is 0.000--0.083 for four of the five defect types, against 0.522--0.960 for PatchCore, and 80 of the 119 anomalous screws are missed only by PaDiM.

\emph{Capsule: background specks, not pose.}
Capsule is as well aligned as bottle (align 0.998), so pose cannot explain PaDiM's deficit.
Instead, 34.7\% of PaDiM's capsule true positives have their map maximum \emph{outside} the GT, against 13.9\% for PatchCore.
Over a near-constant background, the per-position covariance is tiny, so a dust speck becomes a large Mahalanobis outlier that can outscore a subtle defect on the capsule body (Fig.~\ref{fig:failures}b).

\textbf{Localisation gap is smaller.}
Pixel AUROC is dominated by normal pixels and large defects.
Tile has the lowest AU-PRO for both methods (\pdAuproTile{} / \pcAuproTile{}), because its large, low-contrast oil and glue stains are only partly segmented.

\subsection{Threshold calibration}
\label{sec:calib}
PaDiM ranks tile well (image AUROC \pdImgAurocTile{}) and screw moderately, yet its F1 at the validation threshold is only \pdImgFoneTile{} on tile and \pdImgFoneScrew{} on screw.
The corresponding oracle values are \pdImgFmaxTile{} and \pdImgFmaxScrew{}.
Table~\ref{tab:calib} shows why: the threshold sits far above typical normal scores (1.49--1.74$\times$ the validation median, against 1.20 for PatchCore), and on screw it even exceeds the median score of the anomalous test images.
There are two causes.
First, with $d=550$ and only 206--287 degrees of freedom, a \emph{held-out} normal image with energy in the directions the fit data do not span receives an inflated Mahalanobis distance.
This produces a heavy right tail in the validation scores (coefficient of variation 0.15--0.19 vs 0.07 for PatchCore) on the two most variable categories.
We infer this mechanism from the validation distribution; in-sample training distances were not stored.
Second, with only 21--32 validation scores, the interpolated 99\% quantile lies between the two largest scores and is effectively the maximum. On tile, one normal validation image (31.33, against a second-largest score of 24.37) alone sets the threshold at 29.80, above the median anomalous test image.
The result is near-perfect precision at a recall of 0.13--0.17.
PatchCore suffers less but also loses calibration: its mean F1 is \pcImgFoneAll{} against an oracle reference of \pcImgFmaxAll{}.

\begin{table*}[t]
  \centering
  \caption{Calibration on screw and tile. All statistics are image scores; the threshold is the 99\% quantile of the normal validation scores.}
  \label{tab:calib}
  \footnotesize
  \begin{tabular}{@{}l@{\hspace{4pt}}r@{\hspace{4pt}}r@{\hspace{4pt}}r@{\hspace{4pt}}r@{\hspace{4pt}}r@{\hspace{4pt}}r@{\hspace{4pt}}r@{\hspace{4pt}}r@{\hspace{4pt}}r@{}}
\toprule
Method--category & $n_{\rm val}$ & \shortstack{Val\\median} & \shortstack{Val\\max} & \shortstack{Image\\thr} & \shortstack{Thr/val\\median} & \shortstack{Anomaly\\median} & Prec. & Recall & F1 \\
\midrule
PaDiM--screw & 32 & 17.604 & 26.465 & 26.149 & 1.485 & 23.017 & 0.938 & 0.126 & 0.222 \\
PaDiM--tile & 23 & 17.158 & 31.327 & 29.795 & 1.737 & 26.628 & 1.000 & 0.167 & 0.286 \\
PatchCore--screw & 32 & 1.098 & 1.343 & 1.317 & 1.199 & 1.432 & 0.990 & 0.798 & 0.884 \\
PatchCore--tile & 23 & 1.378 & 1.660 & 1.648 & 1.196 & 1.919 & 1.000 & 0.810 & 0.895 \\
\bottomrule
\end{tabular}

\end{table*}
